
\documentclass{article}
\usepackage{graphicx} % Required for inserting images

\usepackage{amsmath}
\usepackage{amsfonts}
\usepackage{tikz-cd}

\title{test}
\date{May 2025}

\begin{document}

\vspace{0.5cm}

\textbf{Exercise 8.3}

\vspace{0.5cm}

Let $L: \ell^2(\mathbb N) \rightarrow \ell^2(\mathbb N):(x_0,x_1,..) \mapsto (x_1,x_2,..)$ be the left shift operator $S = \{L,L^2,L^3,..\}$, then $S$ is closed under operator norm since if $m \neq n$ then $\| L^m -L^n\|_{op} = 2$ so there are no Cauchy sequences w.r.t operator norm (except for the eventually constant ones) contained in $S$ and thus $S$ is closed under operator norm. However $L^m \rightarrow 0$ in the strong operator topology so $0 \in \overline S^{s}$ but $0 \not \in S$ so the strong and norm closures of $S$ differ. 

Now let $R: (x_0,x_1,..) \mapsto (0,x_0,x_1,..)$ be the right shift operator and $S = \{R,R^2,R^3,..\}$. If $(R^{n_i})_{i \in I}$ is a net converging in the strong operator topology to some operator $T$ then let $x = (1,0,0,..)$ then $R^{n_i}x$ has to be a Cauchy net in $\ell^2(\mathbb N)$ so that there exists some $k \in I$ so that $i,j \geq k$ implies $\|R^{n_i}x - R^{n_j}x \| < 1$ but this automatically implies $n_i = n_j$ and thus the net is eventually constant and $T \in S$. However in the weak operator topology $R^n \rightarrow 0$ as $n \rightarrow \infty$ and $0 \not \in S$ so the weak and strong operator closures of $S$ differ.

\vspace{0.5cm}

\textbf{Exercise 8.4}

\vspace{0.5cm}

We need to prove that if the net $P_i \rightarrow P$ in the weak operator topology then $P_i \rightarrow P$ in the strong operator topology.

For an orthogonal projection $Q$ we have 
$$(Qx,x) =(Q^2x,x) = (Qx,Qx) = \|Qx\|^2$$ 
since $Q^2 = Q$ and $Q = Q^*$ so then $\|P_ix\|^2 \rightarrow \|Px\|^2$ by weak operator topology convergence but this together with weak convergence $P_ix \rightharpoonup Px$ implies norm convergence $P_ix \rightarrow Px$ and thus strong operator topology convergence $P_i \rightarrow P$.
\vspace{0.5cm}

\textbf{Exercise 8.5}

\vspace{0.5cm}

Since $\|x\| = \|U_ix\|$ and $\|x\| = \|Ux\|$ then $\|U_ix\| \rightarrow \|Ux\|$ and therefore since also $U_ix \rightharpoonup Ux$ we conclude $U_i x \rightarrow Ux$.

\vspace{0.5cm}

\textbf{Exercise 8.6}

\vspace{0.5cm}

Let's first prove that $\mathcal A = \mathcal K + \mathbb C \text{Id}$ is a $C^*-$algebra, the only unclear part is that it should be closed under operator norm. First let $K,L$ be two compact operators and $z,w \in \mathbb C$, then we would like to establish the inequality 
$$|z-w| \leq \|(K + z \text{Id}) - (L+w\text{Id})\|_{op}$$

Let $e_n$ be an orthonormal basis for $\mathcal H$ then $e_n \rightharpoonup0$ and thus since compact operators take weakly convergent sequences to norm convergent sequences we get $(K-L)e_n \rightarrow 0$ in $\mathcal H$. Thus since
$$\|(z-w)e_n\| - \|(K-L)e_n \| \leq \|(K-L)e_n + (z-w)e_n \| \leq \|(K + z \text{Id}) - (L+w\text{Id})\|_{op}$$
we arrive at the desired inequality after $n \rightarrow \infty$. This inequality we proved implies that if $K_n + \lambda_n \text{Id}$ is Cauchy in operator norm and the $K_n$ are compact that $\lambda_n$ is Cauchy in $\mathbb C$ and therefore that $K_n$ is Cauchy aswell and so 
$$K_n + \lambda_n \text{Id} \rightarrow K + \lambda \text{Id}$$ 
for some compact $K$.

\vspace{0.5cm}

\textbf{Exercise 8.7}

\vspace{0.5cm}

Let $M'$ be the algebra generated by $M$ (all finite combinations through addition and multiplication of elements in $M$) then it is a commutative $*$-algebra with 
$$(M')^\circ = M^\circ$$
we also see that 
$$(M' + \mathbb C \text{Id})^\circ = (M')^\circ = M^\circ$$ 
and thus $M^{\circ \circ} = \overline{M' + \mathbb C \text{Id}}^{SOT}$ by von Neumann's bicommutant theorem. If $x_i \rightarrow x$ and $y_j \rightarrow y$ are nets of operators that converge in the strong operator topology with $x_iy_j = y_j x_i$ then $xy = yx$ thus the SOT closure of $M' + \mathbb C \text{Id}$ is commutative since it's a commutative algebra.

\vspace{0.5cm}

\textbf{Exercise 8.10}

\vspace{0.5cm}

I can only figure this one out with the Borel functional calculus. 

\end{document}
